\begin{myexercise}
    In dieser Übung wollen wir unser Verständnis vom Verhalten von Klassenattributen in der Vererbung stärken.

    \begin{enumerate}
        \item Implementieren Sie die Klasse \texttt{Schulklasse} mit den in Tabelle \ref{tab:attribute_uebersicht} genannten Klassenattributen.
        \item Testen Sie Ihre \textbf{Implementierung} mit folgendem Code:
        \begin{lstlisting}
            print(f"Basis-Handyzeit: {Schulklasse.maximal_handy_zeit}")
            print(f"Basis-Regeln: {Schulklasse.pausenregeln}")
        \end{lstlisting}
        Warum ist es nicht notwendig, dass wir zunächst eine Instanz erstellen?
        \item \textbf{Implementieren} Sie die Klasse \texttt{Klasse1a}. Die Klasse \texttt{Klasse1a} ist eine Kindklasse von \texttt{Schulklasse} und hat eine strengere \texttt{maximal\_handy\_zeit} von 10. 
        \item Testen Sie Ihre \textbf{Implementierung} mit folgendem Code:
        \begin{lstlisting}
            schueler_5a = Klasse5A()
            print(f"1. Instanz-Handyzeit: {schueler_5a.maximal_handy_zeit}")
            print(f"2. Schulklasse Handyzeit: {Schulklasse.maximal_handy_zeit}")
            print(f"3. Instanz-Aufsicht: {schueler_5a.aufsichtsperson}")
        \end{lstlisting}
        \item Erweitern Sie die Klasse \texttt{Klasse5A} um die Methode \texttt{regel\_hinzufuegen()} so, dass es möglich ist, weitere Regeln für diese spezifische Klasse hinzuzufügen.
        \item Testen Sie Ihre \textbf{Implementierung} mit folgendem Code. Gilt die neue Regel wirklich nur für \texttt{Klasse5A}?
        \begin{lstlisting}
            schueler_5a = Klasse5A()
            schueler_5a.regel_hinzufuegen("Keine Handys im Pausenhof") # Aufruf über Instanz oder Klasse

            print(f"1. Instanz Regeln: {schueler_5a.pausenregeln}")
            print(f"2. Schulordnung Regeln: {Schulordnung.pausenregeln}")
        \end{lstlisting}
    \end{enumerate} 
\end{myexercise}

\begin{table}[h]
    \centering
    \caption{Attribute in der Schulklasse}
    \label{tab:attribute_uebersicht}
    \begin{tabular}{|l|l|p{9cm}|l|} 
        \hline
        \textbf{Attributname} & \textbf{Datentyp} & \textbf{Beschreibung} \\
        \hline
        \texttt{maximal\_handy\_zeit} & \texttt{int} & Allgemeine Obergrenze in Minuten (z. B. 15). \\
        \hline
        \texttt{aufsichtsperson} & \texttt{str} & Standard-Aufsichtsperson (z. B. Rektor).  \\
        \hline
        \texttt{pausenregeln} & \texttt{list} & Allgemeine Regeln für die Pause (z. B. \texttt{[Nicht rennen, Müll wegräumen]}). \\
        \hline
    \end{tabular}
\end{table}